% Hexagon note H17, 23 Nov 2026. CC0 1.0.
% Mathematics: GPT-6 Astra (ChatGPT Pro mode), November 2026.
% Text (version 1): Claude Fable 5.2 at effort max, as editor; several proofs rewritten (see "Provenance and checks").
\documentclass[11pt,a4paper]{article}
\usepackage[T1]{fontenc}
\usepackage{lmodern}
\usepackage{microtype}
\usepackage[margin=2.7cm]{geometry}
\usepackage{amsmath,amssymb,amsthm}
\PassOptionsToPackage{hyphens}{url}
\usepackage[hidelinks]{hyperref}
\emergencystretch=2.5em

\theoremstyle{plain}
\newtheorem{theorem}{Theorem}
\newtheorem{lemma}[theorem]{Lemma}
\newtheorem{proposition}[theorem]{Proposition}
\newtheorem{corollary}[theorem]{Corollary}
\theoremstyle{definition}
\newtheorem{definition}[theorem]{Definition}
\theoremstyle{remark}
\newtheorem{remark}[theorem]{Remark}

\newcommand{\GT}{\widehat{\mathrm{GT}}}
\newcommand{\GTz}{\widehat{\mathrm{GT}}_0}
\newcommand{\Zh}{\widehat{\mathbf Z}}
\newcommand{\Z}{\mathbf Z}
\newcommand{\Q}{\mathbf Q}
\newcommand{\F}{\mathbf F}
\newcommand{\GQ}{G_{\mathbf Q}}
\newcommand{\Fh}{\widehat F}
\newcommand{\Bh}{\widehat B}
\newcommand{\Rh}{\widehat R}
\newcommand{\Gh}{\widehat\Gamma}
\newcommand{\Ad}{\operatorname{Ad}}
\newcommand{\Aut}{\operatorname{Aut}}
\newcommand{\Out}{\operatorname{Out}}
\newcommand{\SL}{\operatorname{SL}}
\newcommand{\PSL}{\operatorname{PSL}}
\newcommand{\GL}{\operatorname{GL}}
\newcommand{\per}{\operatorname{per}}
\newcommand{\ab}{\mathrm{ab}}
\newcommand{\recordurl}{https://github.com/KokunoYumeto/mathematics-workbenches/blob/main/works/WB-0007.md}

\title{Relation (IV) of Nakamura and Schneps holds on the profinite Grothendieck--Teichm\"uller group,
with an application to the two elliptic points of the modular curve}
\author{}
\date{23 November 2026 \quad CC0 1.0}

\begin{document}
\maketitle

\begin{abstract}
\noindent
Nakamura and Schneps (2000) defined a subgroup of the profinite Grothendieck--Teichm\"uller group $\GT$ by two
relations, $(\mathrm{III}')$ and $(\mathrm{IV})$, that hold for Galois elements; they proved that it acts on the
profinite mapping class groups of all surfaces, and left open whether it is all of $\GT$. We prove that
$(\mathrm{IV})$ holds for every element of $\GT$. It follows that every element of $\GT$ is the product of a Galois
element and an element of the group $\Lambda$ of Hatcher, Lochak and Schneps, which acts on the whole tower of pure
mapping class groups by their formulas. As an application we study a character $\delta$ with values in
$\mathbf F_2$, defined on the larger group given by the two- and three-cycle relations alone, which compares at
level twelve how an element moves the two elliptic points of the modular curve. We prove that $\delta$ vanishes on
the normalized elements of $\GT$, so that every character extending it and trivial on Galois elements is trivial on
$\GT$. The proofs use the theorem of Boggi and Zalesskii on finite $p$-subgroups of profinite completions, the
action of Hatcher, Lochak and Schneps, and the theorem of Hoshi, Minamide and Mochizuki on the outer automorphisms
of genus-zero configuration groups. Relation $(\mathrm{III}')$ is not touched. The mathematics was produced by
GPT-6 Astra (ChatGPT Pro mode); it was checked by Claude Fable 5.2 at effort max, which wrote this text and rewrote
several proofs, and by Claude Opus 5.5 at effort max as second reader. No human has checked it.
\end{abstract}

\section{The results}

\subsection*{The group and the question}
Let $\Fh_2$ be the free profinite group on $x,y$. For $f\in\Fh_2$ and elements $a,b$ of a profinite group,
$f(a,b)$ is the image of $f$ under $x\mapsto a$, $y\mapsto b$. Drinfeld's group $\GT$ \cite{D} consists of the pairs
$F=(\lambda,f)$, with $\lambda\in\Zh^\times$ and $f$ in the closed commutator subgroup of $\Fh_2$, such that
$x\mapsto x^\lambda$, $y\mapsto f^{-1}y^\lambda f$ is an automorphism of $\Fh_2$ and
\begin{itemize}
\item[(I)] $f(x,y)f(y,x)=1$;
\item[(II)] $f(x,y)x^{m}f(z,x)z^{m}f(y,z)y^{m}=1$, where $xyz=1$ and $m=(\lambda-1)/2$;
\item[(III)] the five-cycle (pentagon) relation in the pure mapping class group $\Gh_{0,5}$.
\end{itemize}
The pairs satisfying (I) and (II) form the larger group $\GTz$ of Harbater and Schneps \cite{HS}, called the gentle
version in \cite{DG}. The absolute Galois group $\GQ$ embeds in $\GT$, with $\lambda$ the cyclotomic character
\cite{D,I}. Whether $\GQ=\GT$ is not known.

Write $B_3=\langle u,v\mid uvu=vuv\rangle$. Every $F\in\GT$ acts on the profinite braid groups $\Bh_n$; on
$\Bh_3$ the action is $u\mapsto u^\lambda$, $v\mapsto f(u^2,v^2)^{-1}v^\lambda f(u^2,v^2)$. Comparing the Galois
actions on $\Gh_{0,5}$ and $\Gh_{1,2}$, Nakamura found that Galois elements satisfy, in $\Bh_3$,
\[
 f(u,v^4)=v^{8\rho_2}\,f(u^2,v^2)\,u^{4\rho_2}\,(uv)^{-6\rho_2},\tag{IV}
\]
where $\rho_2=\rho_2(F)\in\Zh$ is defined for every $F\in\GT$ (Definition~\ref{def:rho}) and is, for a Galois
element, the Kummer cocycle of the roots of $2$ \cite[\S1 and Definition~5.1]{NS}. Nakamura and Schneps defined
$\mathrm{I}\Gamma\subset\GT$ by (IV) and a second relation $(\mathrm{III}')$ in $\Gh_{0,[5]}$, proved that
$\GQ\subset\mathrm{I}\Gamma$ and that $\mathrm{I}\Gamma$ acts on $\Gh(\Sigma)$ for every surface $\Sigma$,
compatibly with inclusions of subsurfaces \cite[Theorems 1.2--1.4]{NS}; they ``leave it open to settle
(in-)equalities between consecutive terms of $\GQ\subset\mathrm{I}\Gamma\subset\GT$'' \cite[\S1]{NS}. Hatcher,
Lochak and Schneps \cite{HLS} had defined the group $\Lambda$ of the $f$ with $(1,f)\in\GT$ that satisfy one
relation (R) in the mapping class group of a torus with two boundary circles, and had shown that $\Lambda$ acts on
the tower of all pure profinite mapping class groups by explicit formulas attached to pants decompositions. In
February 2026 Lochak, Nakamura and Schneps \cite{LNS} proved that all of $\GT$ acts on $\Gh_{g,0}$ and $\Gh_{g,1}$
for every $g\ge1$, and wrote of the added condition that ``it is unclear, however, whether the additional condition
actually defines a proper subgroup of $\GT$'' \cite[\S1]{LNS}.

\begin{theorem}\label{thm:IV}
Relation $(\mathrm{IV})$ holds in $\Bh_3$ for every $F=(\lambda,f)\in\GT$, with $\rho_2=\rho_2(F)$.
\end{theorem}

So the subgroup $\mathrm{I}\Gamma$ of Nakamura and Schneps is defined inside $\GT$ by $(\mathrm{III}')$ alone. For
elements with $\lambda=1$ and $\rho_2=0$ relation (IV) reads $f(u,v^4)=f(u^2,v^2)$, and Schneps showed that this
form implies (R) \cite[p.~670]{S}. Call $F\in\GT$ \emph{normalized} if $\lambda_F=1$ and $\rho_2(F)=0$; the
normalized elements form a closed subgroup $\Lambda_0$ of $\GT$.

\begin{theorem}\label{thm:Lambda}
Every $F\in\GT$ is a product $F=F_*\sigma$ with $\sigma\in\GQ$ and $F_*=(1,f_*)$ normalized, and $f_*\in\Lambda$
for every normalized element. Thus $\GT=\Lambda_0\cdot\GQ$ with $\Lambda_0\subset\Lambda$.
\end{theorem}

Before this, the elements of $\Lambda$ known explicitly were Galois: $\Lambda$ contains the absolute Galois group
of the field generated by all roots of unity and all roots of $2$ \cite[p.~3]{HLS}. Theorem~\ref{thm:Lambda} says
that, up to Galois elements, every element of $\GT$ acts on the whole pure tower by the formulas of \cite{HLS}.

\subsection*{The application}
The modular group $B=B_3/\langle(uv)^3\rangle\cong C_2*C_3$ is the orbifold fundamental group of the moduli of
elliptic curves up to $\pm1$, based at the cusp. Put $s=uvu$, $r=uv$; they generate the stabilizers of the two
elliptic points, the curves $y^2=x^3-x$ and $y^2=x^3+1$ with complex multiplication by $\Z[i]$ and $\Z[\omega]$.
An element $F=(1,f)$ of $\GTz$ acts on $\Bh$ and sends $s$ and $r$ to conjugates
$\beta s\beta^{-1}$ and $ArA^{-1}$ by unique elements $\beta,A$ of the pure subgroup
$\Fh=\overline{\langle u^2,v^2\rangle}$. On a subgroup $H_0$ of index at most three in the elements with
$\lambda=1$, the images of these two \emph{frames} in $\PSL_2(\Z/12)$ define two characters $a_2,a_3:H_0\to\F_2$ and
$\delta=a_2+a_3$ (Section~\ref{sec:delta}). Normalized elements lie in $H_0$.

\begin{theorem}\label{thm:delta}
$\delta(F)=0$ for every normalized $F\in\GT$.
\end{theorem}

\begin{corollary}\label{cor:Delta}
Let $\Delta:\GT\to\F_2$ be a homomorphism that is trivial on $\GQ$ and equal to $\delta$ on $\Lambda_0$. Then
$\Delta=0$.
\end{corollary}

\begin{proof}
By Theorem~\ref{thm:Lambda}, $\Delta(F)=\Delta(F_*)+\Delta(\sigma)=\delta(F_*)=0$.
\end{proof}

\subsection*{What it means, and what is not proved here}
In arithmetic terms the two bits compare two descriptions of one field. For a Galois element with $\lambda=1$ both
record the sign by which it moves $(3/4)^{1/4}$: this number generates the $4$-division field of the curve with
$j=0$ and the field of the $x$-coordinates of the $3$-division points of the curve with $j=1728$, over the
cyclotomic field. This description, the existence of a continuous character $\Delta$ of all of $\GTz$ that
restricts to $\delta$ and vanishes on $\GQ$, and the existence of an element of $\GTz$ on which it is $1$, are
results of the longer text behind this note (see ``Provenance and checks''). \emph{They are cited from that text and
have not been refereed; they are not proved here and are not used in the proofs of Theorems
\ref{thm:IV}--\ref{thm:delta}.} Of the facts behind them, only the equality of the second of the two fields with
$\Q(i,(3/4)^{1/4})$ has been verified independently. Granting them, Corollary~\ref{cor:Delta}
says: the two frames at the elliptic points can disagree on the hexagon group, and the pentagon relation forces
them to agree. The inclusion $\GT\subsetneq\GTz$ itself is known; what is new is an explicit level-twelve character
that separates them, with this meaning.

\subsection*{Credit}
Relation (IV) is Nakamura's, and the group it defines with $(\mathrm{III}')$ is Nakamura and Schneps's; that Galois
elements satisfy it is their theorem \cite[\S1]{NS}. The group $\Lambda$, the relation (R) and the action through
the complex of pants decompositions are due to Hatcher, Lochak and Schneps \cite{HLS}. The implication from the
normalized relation (IV) to (R) is Schneps's \cite[p.~670]{S}. What is added here is the passage from Galois
elements to all of $\GT$ in Theorem~\ref{thm:IV}, its consequence Theorem~\ref{thm:Lambda}, and
Theorem~\ref{thm:delta}. The proofs rest on three theorems of others: the theorem of Boggi and Zalesskii that a
finite $p$-subgroup of the profinite completion of a mapping class group of genus at most two is conjugate into the
discrete group, with the expected normalizer \cite{BZ}; the action of $\Lambda$ \cite{HLS}; and the theorem of
Hoshi, Minamide and Mochizuki that $\Out(\Gh_{0,m})=\GT\times\mathfrak S_m$ for $m\ge5$
\cite[Corollary C]{HMM}.

\subsection*{Context}
The idea of reading relations of $\GT$ off special loci of moduli spaces, where curves have automorphisms, is that
of \cite{S}, where (IV) and $(\mathrm{III}')$ are derived from one special morphism for the elements that respect
it; the action of $\GT$ on torsion of genus-zero mapping class groups was studied by Collas \cite{Co}. Our proof
of Theorem~\ref{thm:IV} is of this kind, with the theorem of \cite{BZ} supplying, for every element of $\GT$, the
conjugation that the earlier works assume. The results of \cite{LNS} concern the groups with at most one marked
point and all of $\GT$, without the compatibility with subsurfaces; they neither use nor imply (IV), and
Theorem~\ref{thm:Lambda} does not give their theorem. Basualdo Bonatto and Robertson \cite{BBR} build a modular
operad of mapping class groupoids on which the subgroup of Nakamura and Schneps acts after completion; by
Theorem~\ref{thm:IV} that subgroup is cut out by $(\mathrm{III}')$ alone. We know of no earlier statement of
Theorems \ref{thm:IV}--\ref{thm:delta}; a search of arXiv and of Hexagon on 23 November 2026 found none.

\subsection*{Scope}
We do not decide whether $(\mathrm{III}')$ holds on $\GT$, whether $\Lambda_0\cdot\GQ$ acts on the tower as one
group (the question of \cite[\S8]{LNS}), or whether $\GQ=\GT$.

\section{Proof overview}

\emph{The difficulty.} Relation (IV) compares $f$ on a half twist and a full twist, $f(u,v^4)$ against
$f(u^2,v^2)$. For Galois elements it expresses that two maps between moduli spaces are defined over $\Q$. An
element of $\GT$ is only an automorphism of profinite groups, and such maps are not part of the structure it is
known to respect.

\emph{The idea.} Use a locus where curves have automorphisms. In the moduli of ten points on the sphere, the
configurations invariant under the Klein four-group $\{\pm t^{\pm1}\}$ form a thrice-punctured sphere. This locus
maps to the moduli of four points in two unordered pairs in two ways, by remembering two different quartets, and the
two maps differ by an automorphism $S$ of the target that exchanges a half-twist cusp with a full-twist cusp. An
element $F$ of $\GT$ acts on $\Gh_{0,[10]}$. By the theorem of Boggi and Zalesskii the image of the four-group is
conjugate to a discrete four-group, hence, after a correction that both forgetting maps kill, equal to it; $F$ then
acts on its centralizer, compatibly with both maps. So the standard action on the four-point group commutes with
$S$ up to an inner automorphism (Proposition~\ref{prop:swap}). Centralizers in a free product turn this into one
equation, and abelianization on a subgroup of index two gives (IV) with a coefficient $c$; a second character
computation shows $c=\rho_2$ (Section~\ref{sec:IV}).

\emph{From (IV) to the tower.} Multiplying by a Galois element makes $\lambda=1$ and $\rho_2=0$; Schneps's
computation gives (R), so the action of \cite{HLS} is available (Section~\ref{sec:Lambda}). We extend it to mapping
classes that permute marked points and check that it commutes with forgetting a point (Section~\ref{sec:genusone}).

\emph{The character.} For Theorem~\ref{thm:delta} the same method is used in genus one, twice. The translations by
points of order three form a group $C_3^2$ in the mapping class group of a torus with nine marked points; its
normalizer contains a copy of $\SL_2(\Z)$, on which a normalized $F$ acts as on $\Gh_{1,1}$. Dividing by the
elliptic involution and forgetting three branch points leads to a sphere with five marked points, and the double
cover branched at the images of the points of order three is again a torus with two marked points. A classical
identity makes this torus the curve $y^2=x^3+1$, whatever the curve one starts from: the modular orbifold is
identified with the quotient of the punctured hexagonal curve by its automorphism group $\mu_6$. So $F$ acts
compatibly on the modular group and on a subgroup $G_3$ of $\Gh_{1,2}$ isomorphic to it
(Proposition~\ref{prop:H}). Under this isomorphism the frame of $F$ at the point of order three becomes a
six-fold symmetry of a lattice, the frame at the point of order two a translation, and the two bits are read off
the same matrix modulo $4$ and modulo $3$ (Section~\ref{sec:proof}). The hypothesis $\rho_2=0$ is used to have the
action of \cite{HLS}; the pentagon is used through \cite{HLS} and \cite{HMM}.

\section{Conventions}\label{sec:conv}

\subsection*{Groups}
Completions are profinite completions; $\overline{X}$ is the closure of $X$. We multiply in $\GTz$ so that
$F\mapsto\Phi_F$ is a homomorphism, where $\Phi_F(x)=x^\lambda$, $\Phi_F(y)=f^{-1}y^\lambda f$. For
$\sigma\in\GQ$, $\lambda$ is the cyclotomic character. $\Gamma_{g,[n]}$ is the mapping class group of an oriented
surface of genus $g$ with $n$ marked points, which may be permuted; $\Gamma_{g,n}$ is its pure subgroup. A subscript
such as $[9]+[3]$ means that two sets of marked points are preserved separately.

\subsection*{Braids}
$B_n$ has generators $\sigma_1,\dots,\sigma_{n-1}$; in $B_3$ we write $u,v$, and in $B_4$ we write $a,b,c$. Put
$y_i=\sigma_{i-1}\cdots\sigma_1^2\cdots\sigma_{i-1}$. We use the following known facts.
\begin{itemize}
\item[(B1)] For $F\in\GTz$ the formulas $u\mapsto u^\lambda$, $v\mapsto f(u^2,v^2)^{-1}v^\lambda f(u^2,v^2)$ define an
automorphism $\Phi_F$ of $\Bh_3$, with $\Phi_F((uv)^3)=(uv)^{3\lambda}$ \cite{LS,HS}.
\item[(B2)] For $F\in\GT$ and $n\ge3$ the formulas $\sigma_1\mapsto\sigma_1^\lambda$,
$\sigma_i\mapsto f(y_i,\sigma_i^2)^{-1}\sigma_i^\lambda f(y_i,\sigma_i^2)$ define an automorphism $A_F$ of $\Bh_n$
\cite{D,LS}, \cite[Theorem 2.1]{LNS}. It induces the identity on permutations and descends to the quotient
$\Gh_{0,[n]}$ and to the quotient by the centre, which is the stabilizer of one point in $\Gh_{0,[n+1]}$.
\item[(B3)] \emph{Forgetting.} Let $M_S\subset\Gamma_{0,[n]}$ be the stabilizer of a set $S$ of marked points and
$p_S:M_S\to\Gamma_{0,[S]}$ the map forgetting the others. For $F\in\GT$ there is a pure $g\in\Gh_{0,S}$ with
$p_S\circ A_F=\Ad(g)\circ A_F\circ p_S$ on $\widehat M_S$.
\end{itemize}
For (B3), see \cite[\S3.1]{HS}; the proof is short. If $S$ is the set of the first $k$ strands, $p_S$ is induced by
deleting the last $n-k$ strands of braids that preserve the first $k$; this group is generated by the $\sigma_i$
with $i\ne k$ and by $\sigma_k^2$. For $i<k$ the formula for $A_F(\sigma_i)$ is the same word in $\Bh_k$; for $i>k$,
and for $\sigma_k^2$, the generator dies and its image under $A_F$, a conjugate of a power of it by a pure braid,
dies too. So $p_S$ commutes with $A_F$. In general conjugate $S$ to the first $k$ strands by a braid $b$, and note
that $A_F(b)b^{-1}$ is pure.

\begin{itemize}
\item[(F)] \emph{Completion of a free normal subgroup.} Let $\Pi$ be a normal subgroup of a group $A$ that is free
of finite rank $r\ge2$. Then the profinite topology of $A$ induces the full profinite topology on $\Pi$, and
$1\to\widehat\Pi\to\widehat A\to\widehat{A/\Pi}\to1$ is exact; see also \cite{A}.
\end{itemize}
For (F), let $N$ be a characteristic subgroup of finite index of $\Pi$, $D=\Pi/N$, $\ell$ a prime larger than
$|D|(r-1)$, and $N_1=[N,N]N^\ell$. On the cover with group $D$ of a rose with $r$ petals, a nontrivial element of
$D$ has trace $0$ on the cellular chains, hence trace $1$ on $N/N_1=H_1(N,\F_\ell)$, whose dimension
$1+|D|(r-1)$ is not $1$ modulo $\ell$. So $D$ acts faithfully on $N/N_1$, the centre of $\Pi/N_1$ has trivial image
in $D$, and the kernel of the finite quotient $A\to\Aut(\Pi/N_1)$ meets $\Pi$ inside $N$.

\subsection*{The modular group, frames and $\rho_2$}
Let $B=B_3/\langle(uv)^3\rangle$, $s=uvu$, $r=uv$, $\tau=sr$. Then $s^2=r^3=1$, $B=\langle s\rangle*\langle
r\rangle$ and $\tau=u^{-1}$. The elements $x=u^2$, $y=v^2$ freely generate the kernel $F$ of $B\to\mathfrak S_3$,
and $sxs^{-1}=y$. We identify $\Fh$ with $\Fh_2$; by (B1) every $F\in\GTz$ acts on
$\Bh=\widehat{C_2}\amalg\widehat{C_3}$, by the identity on $\mathfrak S_3$ and by $\Phi_F$ on $\Fh$. Let
$\vartheta=\Ad(s)|_{\Fh}$, which exchanges $x$ and $y$.

We use the theorem of Herfort and Ribes \cite{HR}: in a free profinite product $G_1\amalg G_2$, every finite
subgroup is conjugate into a factor, and the centralizer of a nontrivial element of a factor lies in that factor.
In particular $C_{\Bh}(s)=\langle s\rangle$, $C_{\Bh}(r)=\langle r\rangle$, and
$C_{\Fh}(y)=\overline{\langle y\rangle}$.

\begin{lemma}\label{lem:frames}
Let $F\in\GTz$. There is a unique $\beta_F\in\Fh$ with $\Phi_F(s)=\beta_Fs\beta_F^{-1}$. If $\lambda_F=1$ there is
a unique $A_F\in\Fh$ with $\Phi_F(r)=A_FrA_F^{-1}$. Moreover $\beta_{FG}=\Phi_F(\beta_G)\beta_F$, and
$A_{FG}=\Phi_F(A_G)A_F$ when $\lambda_F=\lambda_G=1$.
\end{lemma}

\begin{proof}
$\Phi_F(s)$ has order two, so it is $gsg^{-1}$ for some $g\in\Bh$. Since $\Phi_F$ is the identity on $\mathfrak
S_3$, the image of $g$ commutes with that of $s$, so $g$ or $gs$ lies in $\Fh$. Uniqueness holds because
$C_{\Bh}(s)\cap\Fh=1$. If $\lambda=1$, $\Phi_F$ is the identity on $B^{\ab}=\Z/6$, where $r\ne r^{-1}$; so the
element $\Phi_F(r)$ of order three is conjugate to $r$, and the same argument applies. The product formulas follow
from uniqueness.
\end{proof}

\begin{definition}\label{def:rho}
For $F\in\GTz$, $\rho_2(F)\in\Zh$ is defined by $[\beta_F^{-1}]=(\rho_2(F),\rho_2(F))$ in
$\Fh^{\ab}=\Zh x\oplus\Zh y$.
\end{definition}

\begin{lemma}\label{lem:rho}
\textup{(a)} $f=\vartheta(\beta_F)\beta_F^{-1}$; so $g_F=\beta_F^{-1}$ satisfies $f(x,y)=g_F(y,x)^{-1}g_F(x,y)$, the
class of $\beta_F^{-1}$ is of the form $(\rho,\rho)$, and $\rho_2$ is the coefficient of \cite[Definition 5.1]{NS}.
\textup{(b)} $\rho_2(FG)=\rho_2(F)+\lambda_F\rho_2(G)$.
\textup{(c)} If $\lambda_F=1$ then $\rho_2(F)\in2\Zh$.
\textup{(d)} For $\sigma\in\GQ$, $\rho_2(\sigma)$ is the Kummer cocycle $\kappa_2$ of the positive roots of $2$:
$\sigma(2^{1/n})=\zeta_n^{\kappa_2(\sigma)}2^{1/n}$.
\end{lemma}

\begin{proof}
(a) From $y=sxs^{-1}$ we get
$\Phi_F(y)=\beta s\beta^{-1}x^\lambda\beta s^{-1}\beta^{-1}
=\beta\vartheta(\beta)^{-1}y^\lambda\vartheta(\beta)\beta^{-1}$. Comparing with $f^{-1}y^\lambda f$, the element
$f\beta\vartheta(\beta)^{-1}$ centralizes $y^\lambda$, so it is $y^e$ for some $e\in\Zh$. Abelianizing, with
$[\beta]=(b_1,b_2)$ and $[f]=0$, gives $(b_1-b_2,b_2-b_1)=(0,e)$; so $b_1=b_2$ and $e=0$.
(b) $\Phi_F$ acts on $\Fh^{\ab}$ by $\lambda_F$; apply Lemma~\ref{lem:frames}.
(c) In the quotient of $\Fh$ of nilpotency class two, relation (II) with $\lambda=1$ gives $[x,y]^{3e}=1$ for
$f=[x,y]^e$, so $f=1$ there. Writing $g_F=x^\rho y^\rho[x,y]^d$ one finds
$g_F(y,x)^{-1}g_F(x,y)=[x,y]^{\rho^2+2d}$, so $\rho^2\in2\Zh$.
(d) This is \cite[\S5, after Definition 5.1]{NS}.
\end{proof}

\section{Relation (IV)}\label{sec:IV}

\subsection*{The coloured four-point group}
Let $R_0=\langle d,z\mid(zd)^2=1\rangle$ and $R=\widehat{R_0}$. With $a=zd$ we have
$R_0=\langle a\rangle*\langle d\rangle\cong C_2*\Z$. The map $\iota_0:d\mapsto u$, $z\mapsto v^2$ identifies $R_0$
with the subgroup of index three of $B$ that is the image of $J=\langle u,v^2\rangle\subset B_3$; one has
$(uv^2)^2=(uv)^3$, and $J$ has the presentation $\langle u,z\mid(uz)^2=(zu)^2\rangle$. Geometrically $R_0$ is the
orbifold fundamental group of the moduli $\mathcal R$ of four points on the sphere in two unordered pairs: for
the quartet $\{0,\infty\},\{t,1/t\}$ the coordinate is $w=t^2+t^{-2}$, and $\mathcal R$ is the $w$-line with a cusp
at $w=\infty$, where points of different pairs meet, a cusp at $w=2$, where the two points of a pair meet, and a
point of order two at $w=-2$. The loops $z$ and $d$ go around $\infty$ and $2$. In terms of mapping classes, $R_0$
is the group of classes of the four-pointed sphere preserving each pair, modulo the class that flips both pairs;
it is the image of that group in $\Gamma_{0,[4]}/V=B$, where $V$ is the normal four-group.

By (B1), $F\in\GTz$ acts on $R$ through $\iota_0$:
\[
 \Phi_R(d)=d^\lambda,\qquad \Phi_R(z)=h\,z^\lambda h^{-1},\qquad h=f(z,d^2).
\]
For $F\in\GT$ this is the action induced by $A_F$ on $\Gh_{0,[4]}$, because the projection of the sphere braid
group on four strands to $B$ sends $(\sigma_1,\sigma_2,\sigma_3)$ to $(u,v,u)$ and $f(vu^2v,u^2)=1$ in $\Bh$. Let
$S$ be the automorphism of $R_0$ exchanging $d$ and $z$.

\begin{proposition}\label{prop:swap}
For every $F\in\GT$ there is $g_0\in R$ with $S\Phi_RS=\Ad(g_0)\Phi_R$.
\end{proposition}

The proof occupies the next three subsections.

\subsection*{Finite $p$-subgroups}
\begin{lemma}\label{lem:BZ}
Let $\Gamma$ be a subgroup of finite index of a mapping class group $\Gamma_{g,[n]}$ with $g\le2$ and $2g-2+n>0$,
let $P$ be a normal subgroup of finite index of $\Gamma$, and let $H\subset\Gamma$ be a finite $p$-group with
$H\cap P=1$.
\begin{itemize}
\item[(a)] Let $\alpha$ be a continuous automorphism of $\widehat\Gamma$ that preserves $\widehat P$ and induces the
identity on $\Gamma/P$. Then there are $e\in\widehat P$ and a homomorphism $h\mapsto h'$ from $H$ to $\Gamma$ with
$h'\in Ph$ and $\alpha(h)=eh'e^{-1}$ for all $h\in H$.
\item[(b)] $N_{\widehat\Gamma}(H)=\overline{N_\Gamma(H)}$ and $C_{\widehat\Gamma}(H)=\overline{C_\Gamma(H)}$.
\end{itemize}
The same holds for a group $\Gamma=K\rtimes H$ when $K$ has a subgroup of finite index that is an iterated
extension of finitely generated free groups.
\end{lemma}

\begin{proof}
Mapping class groups of genus at most two, and hence their subgroups of finite index, are good and of finite
virtual cohomological type, so the theorem of Boggi and Zalesskii applies to them: every finite $p$-subgroup of the
completion is conjugate to a subgroup of the discrete group, and normalizers of finite $p$-subgroups of the
discrete group are the closures of the discrete normalizers \cite[Theorem A and the theorem on hyperelliptic
mapping class groups of the last section]{BZ}. The groups $K\rtimes H$ of the last sentence satisfy the same
hypotheses. (a) Apply this to the group $PH=P\rtimes H$, whose completion is $\widehat P\rtimes H$ and contains
$\alpha(H)$: there are $g=e h_0$, with $e\in\widehat P$ and $h_0\in H$, and a finite $H''\subset PH$ with
$\alpha(H)=gH''g^{-1}$. Put $h'=e^{-1}\alpha(h)e\in h_0H''h_0^{-1}\subset\Gamma$; modulo $\widehat P$ it is $h$.
(b) The first equality is the cited statement. The centralizer is the kernel of
$N_{\widehat\Gamma}(H)\to\Aut(H)$, hence open in the closure of $N_\Gamma(H)$; an open subgroup of the closure of a
subgroup $N$ is the closure of its intersection with $N$, which here is $C_\Gamma(H)$.
\end{proof}

\subsection*{The Klein configuration}
Let $H=\{1,a,b,ab\}$ act on $\mathbf P^1$ by $a(t)=-t$, $b(t)=1/t$, and let
$Z=\mathbf P^1\setminus\{0,\infty,\pm1,\pm i\}$. For $t\in Z$ let $x(t)$ be the configuration of ten labelled
points
\[
 A_\pm=(0,\infty),\quad B_\pm=(1,-1),\quad C_\pm=(i,-i),\quad D_h=h(t)\ (h\in H).
\]
It is invariant under $H$, which permutes the labels. Fix $t_0=1+\sqrt2$ and use $x(t_0)$ as reference surface.
Then $H$ is a subgroup of $M=\Gamma_{0,[10]}$. Let $G\subset M$ be the group of classes whose permutation of the
labels commutes with those of $H$, $P=\Gamma_{0,10}$, and $N=C_G(H)$.

\emph{The discrete centralizer.} A class in $N$ is represented by a homeomorphism commuting with $H$, because the
fixed locus of $H$ in Teichm\"uller space is connected \cite{K}. It descends to the quotient sphere
$\mathbf P^1/H$, with coordinate $w=t^2+t^{-2}$, and fixes there the images $\infty,2,-2$ of the pairs $A,B,C$,
whose stabilizers in $H$ are different, and the image of $D$. Hence $N/H$ is the pure mapping class group of a
sphere with four marked points, that is,
\[
 N/H\;\cong\;\pi_1\bigl(\mathbf P^1_w\setminus\{\infty,2,-2\},\,w_0\bigr),\qquad w_0=6,
\]
a loop being the motion of the point $D$.

\emph{The correction.} Let $F\in\GT$ and $\alpha=A_F$ on $\widehat M$. By Lemma~\ref{lem:BZ}(a), applied to
$\Gamma=G$, there are $e\in\widehat P$ and $h\mapsto h'$ with $\alpha(h)=eh'e^{-1}$, where $h'$ and $h$ permute
the labels in the same way. The finite group $\{h'\}$ acts by M\"obius transformations on a sphere with ten marked
points \cite{K}. Put $A_+,A_-,B_+$ at $0,\infty,1$. Then $a'$, which fixes $A_\pm$ and has order two, is $t\mapsto-t$;
$b'$ exchanges $0,\infty$ and fixes $1$, so it is $t\mapsto1/t$; hence $B_-=-1$, the fixed points $C_\pm$ of $a'b'$
are $\pm i$ in one of two orders, and $D_h=h(t_1)$ for some $t_1\in Z$. The configuration is $x(t_1)$, or $x(t_1)$
with the labels $C_+,C_-$ exchanged. Since $Z$ is connected, a path from $t_0$ to $t_1$ gives $p\in P$ with
$h'=php^{-1}$ in the first case and $h'=p\theta h\theta^{-1}p^{-1}$ in the second, where $\theta\in G$ is a half
twist exchanging $C_+$ and $C_-$ in a disc containing no other marked point. So there is
$e_1\in\widehat P\cup\widehat P\theta$ such that
\[
 \alpha'=\Ad(e_1^{-1})\circ A_F\quad\text{fixes $H$ pointwise.}
\]
By Lemma~\ref{lem:BZ}(b), $\alpha'$ preserves $C_{\widehat G}(H)=\overline N$ and induces an automorphism
$\alpha_N$ of $\overline N/H$.

\subsection*{The two forgetting maps}
Let $M_b\subset M$ be the stabilizer of the sets $\{A_+,A_-\}$ and $\{D_1,D_b\}$. Then $N_b=N\cap M_b$ has index two
in $N$, $N=HN_b$ and $H\cap N_b=\langle b\rangle$. Forgetting the six other points maps $M_b$ to the group of
classes of a four-pointed sphere preserving two pairs, and $b$ to the class flipping both. So we get
$q_b:N/H\to R_0$. On the locus $Z$ the retained quartet is $\{0,\infty\},\{t,1/t\}$, with coordinate $w$. Hence
$q_b$ is induced by the inclusion of $\mathbf P^1_w\setminus\{\infty,2,-2\}$ in $\mathcal R$; it is surjective.
In the same way the stabilizer $M_a$ of $\{B_+,B_-\}$ and $\{D_1,D_a\}$ gives $q_a:N/H\to R_0$. Its quartet
$\{1,-1\},\{t,-t\}$ is carried by the map $\xi\mapsto(\xi-1)/(\xi+1)$ to $\{0,\infty\},\{r,1/r\}$ with
$r=(t-1)/(t+1)$, and
\[
 r^2+r^{-2}=\frac{2(w+6)}{w-2}=:T(w).
\]
$T$ is an involution of $\mathcal R$ that fixes $w=-2$ and $w_0=6$ and exchanges the two cusps. So
$q_a=T_*\circ q_b$, where $T_*$ maps the conjugacy class of $d$ to that of $z$ and conversely. We claim
that $T_*=\Ad(g_1)S$ with $g_1\in R_0$. By the generators of Fouxe-Rabinovitch for the automorphism group of a
free product \cite{FR}, $\Aut(R_0)$ is generated by the inner automorphisms, the partial conjugations, which here
are inner, and the automorphisms $\iota:d\mapsto d^{-1}$, $t:d\mapsto da$ and $t':d\mapsto ad$, each fixing $a$.
The order of $\Out(R_0)$ is not stated in \cite{FR}; we derive it. One has $t'=\Ad(a)t$, $t^2=1$ and
$\iota t\iota=\Ad(a)t$, so $\Out(R_0)$ consists of the classes of $1,\iota,t,t\iota$. They send the image $(1,0)$
of $d$ in $R_0^{\ab}=\Z\times\Z/2$ to $(1,0)$, $(-1,0)$, $(1,1)$, $(-1,1)$, so they are distinct. The image of
$z=ad^{-1}$ is $(-1,1)$. Hence a class that exchanges the conjugacy classes of $d$ and $z$ is that of $t\iota$;
both $T_*$ and $S$ are such classes.

\begin{proof}[Proof of Proposition~\ref{prop:swap}]
Both maps extend to $\overline N/H\to R$, surjectively. The element $e_1$ lies in the completion of $M_b\cap M_a$,
and its images under the two forgetting maps are pure, because $\theta$ is supported in a disc that contains only
forgotten points. By (B3) and the description of $\Phi_R$,
\[
 q_b\alpha_N=\Ad(g_b)\Phi_Rq_b,\qquad q_a\alpha_N=\Ad(g_a)\Phi_Rq_a
\]
for some $g_a,g_b\in R$. Substituting $q_a=\Ad(g_1)Sq_b$ in the second equation and using the first gives
$\Ad(g_1)S\Ad(g_b)\Phi_Rq_b=\Ad(g_a)\Phi_R\Ad(g_1)Sq_b$. As $q_b$ is surjective, $S\Phi_R$ and $\Phi_RS$ differ by
an inner automorphism.
\end{proof}

\subsection*{From the proposition to the relation}
Let $E\subset\Fh_2$ be the kernel of $x\mapsto1$, $y\mapsto0$ in $\Z/2$. It is free on $x^2$, $y$, $xyx^{-1}$, and
$f$, which has exponent sums zero, has class $(0,c,-c)$ in $E^{\ab}$ for a unique $c=c(f)\in\Zh$.

\begin{proposition}\label{prop:IVc}
For every $F\in\GT$: $f(u,v^4)=v^{8c}f(u^2,v^2)u^{4c}(uv)^{-6c}$ in $\Bh_3$, with $c=c(f)$.
\end{proposition}

\begin{proof}
Put $k=S(h)=f(d,z^2)$. Then $S\Phi_RS(z)=z^\lambda$ and $S\Phi_RS(d)=kd^\lambda k^{-1}$. Evaluate
$S\Phi_RS=\Ad(g_0)\Phi_R$ on $z$ and on $d$: $g_0h$ centralizes $z^\lambda$ and $k^{-1}g_0$ centralizes
$d^\lambda$. By \cite{HR}, applied to $R=\langle a\rangle\amalg\overline{\langle d\rangle}$ and to its image under
$S$, these centralizers are $\overline{\langle z\rangle}$ and $\overline{\langle d\rangle}$. So
$g_0=z^sh^{-1}=kd^t$ with $s,t\in\Zh$, that is,
\[
 z^s=k\,d^t\,h.
\]
The kernel $K_0$ of $R_0\to\Z/2$, $d\mapsto1$, $z\mapsto0$, is free on $z$ and $d^2$, since $dzd^{-1}=z^{-1}d^{-2}$.
The elements $h=f(z,d^2)$ and $k$ lie in $\widehat{K_0}$, so $t\in2\Zh$. In $K_0^{\ab}$, with basis $z,d^2$, the
class of $h$ is $0$. The substitution $x\mapsto d$, $y\mapsto z^2$ sends $x^2,y,xyx^{-1}$ to
$d^2,z^2,(z^{-1}d^{-2})^2$, of classes $(0,1),(2,0),(-2,-2)$; so the class of $k$ is $c(2,0)-c(-2,-2)=(4c,2c)$.
Abelianizing the displayed equation gives $(s,0)=(4c,2c+t/2)$: $s=4c$, $t=-4c$, and
$f(d,z^2)=z^{4c}f(d^2,z)d^{4c}$ in $R$. The preimage of $K_0$ in $J$ is $K_0\times\langle(uv)^3\rangle$, so the
kernel of $\widehat J\to R$ is $\overline{\langle(uv)^3\rangle}\cong\Zh$, and
$f(u,v^4)=v^{8c}f(u^2,v^2)u^{4c}(uv)^{3e}$ for some $e\in\Zh$. The character $u\mapsto1$, $v^2\mapsto0$ of $J$ takes
the value $2$ on $(uv)^3=(uv^2)^2$; applied to this equation it gives $0=4c+2e$.
\end{proof}

\begin{proposition}\label{prop:c}
$c(f)=\rho_2(F)$ for every $F\in\GT$.
\end{proposition}

\begin{proof}
Let $j=\Ad(u^{-1}s)\circ\iota_0:R\to\Bh$. Then $j(z)=x$, $j(d)=x^{-1}s$, $j(a)=s$, $j(d^2)=x^{-1}y^{-1}$, and the
image of $j$ contains $\Fh$. Let $\varphi=\Ad(h^{-1})\Phi_R$, so $\varphi(z)=z^\lambda$,
$\varphi(d)=h^{-1}d^\lambda h$. Since $\iota_0\Phi_R=\Phi_F\iota_0$, we have
$j\Phi_Rj^{-1}=\Ad(c_0\Phi_F(c_0)^{-1})\Phi_F$ with $c_0=u^{-1}s$, and
$c_0\Phi_F(c_0)^{-1}=u^{-1}s\beta s^{-1}\beta^{-1}u^{\lambda}=u^{-1}fu^\lambda$ by Lemma~\ref{lem:rho}(a). Also
$j(h)=f(x,x^{-1}y^{-1})=u^{-1}fu$, because $u^{-1}yu=x^{-1}y^{-1}$ in $B$. Hence
\[
 j\varphi j^{-1}=\Ad(u^{\lambda-1})\circ\Phi_F .
\]
Applying this to $j(a)=s$ gives $\varphi(a)=bab^{-1}$ with $b=z^{m}j^{-1}(\beta_F)$, $m=(\lambda-1)/2$.
From the proof of Proposition~\ref{prop:IVc}, $S\Phi_RS=\Ad(kd^{-4c})\Phi_R$; as $\Phi_R=\Ad(h)\varphi$ and
$S(h)=k$, this is $S\varphi S=\Ad(g)\varphi$ with $g=d^{-4c}h$. Evaluate both sides on $a$, using $S(a)=dad^{-1}$
and $S(\varphi(d))=k^{-1}z^\lambda k$: the left side is the conjugate of $a$ by
$k^{-1}z^\lambda k\cdot S(b)\cdot d$, the right side its conjugate by $gb$. The quotient of the two conjugators
centralizes $a$, so it lies in $\langle a\rangle$ \cite{HR}. Apply the character $L:R\to\Zh$ with $L(z)=1$,
$L(d)=-1$, which kills $a$ and satisfies $L\circ S=-L$. Since $L(j^{-1}x)=L(j^{-1}y)=1$, we have
$L(b)=m-2\rho_2$, and $L(h)=L(k)=0$. Therefore
\[
 4c+m-2\rho_2=L(gb)=\lambda-(m-2\rho_2)-1,
\]
that is, $4(c-\rho_2)=\lambda-1-2m=0$.
\end{proof}

Propositions \ref{prop:IVc} and \ref{prop:c} prove Theorem~\ref{thm:IV}. In the notation of \cite[(5.2)]{NS},
where $f\equiv A_F\cdot[x,y]$ modulo the second derived group, one has $c(f)=-A_F(-1,1)$: on $E^{\ab}$ conjugation
by $x$ exchanges the last two coordinates, $y$ acts trivially, and $[x,y]$ has class $(0,-1,1)$. So
Proposition~\ref{prop:c} is the equality $\rho_2=-A_F(-1,1)$ on all of $\GT$, which Nakamura and Schneps prove
under $(\mathrm{III}')$ \cite[Proposition 5.3]{NS}.

\section{Normalization and the group \texorpdfstring{$\Lambda$}{Lambda}}\label{sec:Lambda}

\begin{proof}[Proof of Theorem~\ref{thm:Lambda}]
Let $F\in\GT$. Choose $\sigma_0\in\GQ$ with cyclotomic character $\lambda_F$; then $F_0=F\sigma_0^{-1}$ has
$\lambda=1$, and $\rho_2(F_0)\in2\Zh$ by Lemma~\ref{lem:rho}(c). On the Galois group of $\Q^{\ab}$ the Kummer
cocycle $\kappa_2$ is a homomorphism with image $2\Zh$: the image is a closed subgroup, it lies in $n\Zh$ exactly
when $2^{1/n}\in\Q^{\ab}$, and this happens for $n=1,2$ only, because for $n>2$ the field $\Q(2^{1/n})$ is not
normal over $\Q$. So there is $\sigma_1$ fixing $\Q^{\ab}$ with $\kappa_2(\sigma_1)=\rho_2(F_0)$, and
$F_*=F_0\sigma_1^{-1}$ is normalized by Lemma~\ref{lem:rho}(b),(d).

Let $F=(1,f)$ be normalized. By Theorem~\ref{thm:IV}, $f(u,v^4)=f(u^2,v^2)$ in $\Bh_3$; exchanging $u$ and $v$ and
using (I), also $f(u^4,v)=f(u^2,v^2)$. Relation (R) of \cite[p.~4]{HLS} is the equation $\Pi=1$ in $\Gh_1^2$, the
completed mapping class group of the torus with two boundary circles, where
\[
 \Pi=f(e_3,a)f(b^2,c^2)f(e_2,e_3)f(e_1,e_2)f(a^2,b^2)f(c,e_1),
\]
$a,b,c$ are the twists along a chain of three curves, $e_1=(ab)^6$, $e_3=(bc)^6$ and
$e_2=cbab^{-1}c^{-1}$ \cite[p.~659]{S}. Read $\Pi$ in $\Bh_4$. Schneps's computation \cite[p.~670]{S} shows that
its image in $\Bh_4$ modulo the centre is trivial: the two forms of (IV) just written turn the first, third,
fourth and sixth factors into values of $f$ on pure braids, the two middle factors combine by (II), and the result
is (III). So $\Pi$ is a power $\Omega^n$ of the central element $\Omega=(abc)^4$. But $\Pi$ lies in the closed
commutator subgroup, and $\Omega$ has degree $12$ in $\Bh_4^{\ab}=\Zh$; so $n=0$ and $\Pi=1$ in $\Bh_4$. The
homomorphism $\Bh_4\to\Gh_1^2$ sending $a,b,c$ to the three twists gives (R). Hence $f\in\Lambda$.
\end{proof}

\section{The character \texorpdfstring{$\delta$}{delta}}\label{sec:delta}

\subsection*{Level twelve}
Let $U=\left(\begin{smallmatrix}1&1\\0&1\end{smallmatrix}\right)$,
$L=\left(\begin{smallmatrix}1&0\\-1&1\end{smallmatrix}\right)$, $S_0=ULU=\left(\begin{smallmatrix}0&1\\-1&0
\end{smallmatrix}\right)$, $R_+=UL$ and $R_0=-R_+$; so $R_+^3=-I$. For a ring $A$, $u\mapsto U$, $v\mapsto L$
defines $\tilde q_A:\Bh_3\to\SL_2(A)$ when $\SL_2(A)$ is profinite. We write $\tilde q$ for $A=\Z/12$ and
$q:\Bh\to\PSL_2(\Z/12)=\SL_2(\Z/12)/\{\pm I\}$ for the induced map. On the free group $\Fh$ let
$q_{12}:\Fh\to\SL_2(\Z/12)$ be $x\mapsto U^2$, $y\mapsto L^2$, and $E_{12}$ its image.

\begin{lemma}\label{lem:E}
\textup{(a)} Under $\Z/12=\Z/4\times\Z/3$, $E_{12}=V_4\times\SL_2(\F_3)$, where $V_4=\{I,U^2,L^2,U^2L^2\}$ modulo
$4$ is a four-group; $E_{12}$ does not contain $-I$, so it maps injectively to $\PSL_2(\Z/12)$.
\textup{(b)} The elements of $E_{12}$ whose class commutes with that of $R_+$ form
$\{I\}\times(\langle-I\rangle\times\langle R_0\rangle)$.
\textup{(c)} Those whose class commutes with that of $S_0$ form $\{I,U^2L^2\}\times\langle S_0\rangle$.
\end{lemma}

\begin{proof}
(a) Modulo $4$ the matrices $U^2,L^2$ commute and have order two; modulo $3$ they are $U^{-1},L^{-1}$, which
generate $\SL_2(\F_3)$. The cube of $(U^2,U^{-1})$ is $(U^2,I)$, so $E_{12}$ is the full product. No element of
$V_4$ is $-I$. (b) Let $(g_4,g_3)$ conjugate $R_+$ to $\pm R_+$. The sign is $+$, because $R_+$ has trace $1$,
which is not $-1$ modulo $4$. Conjugation by $R_+$ permutes the three nontrivial elements of $V_4$ cyclically up
to sign, so $g_4=I$; and the centralizer of $R_+$ in $\SL_2(\F_3)$ is $\langle R_+\rangle$. (c) Conjugation by
$S_0$ exchanges $U^2$ and $L^2$, so $g_4S_0g_4^{-1}$ is $S_0$ or $U^2L^2S_0$; the latter is not $-S_0$. So the
sign is $+$, $g_4\in\{I,U^2L^2\}$, and $g_3$ lies in the centralizer $\langle S_0\rangle$ of $S_0$ in
$\SL_2(\F_3)$.
\end{proof}

\begin{lemma}\label{lem:shear}
Let $F\in\GTz$ with $\lambda_F=1$. There is a unique $t_F\in\Z/12$ with
$\tilde q\circ\Phi_F=\Ad(U^{t_F})\circ\tilde q$ on $\Bh_3$. The map $F\mapsto t_F$ is a continuous homomorphism,
$t_F\equiv0\bmod4$, and $t_F\equiv-\rho_2(F)\bmod3$.
\end{lemma}

\begin{proof}
Let $\ell\in\{2,3\}$ and $M=\tilde q_{\Z_\ell}(\Phi_F(v))$. Then $M$ is conjugate to $L$, so it has trace $2$, and
$UMU=MUM$. Writing $M=\left(\begin{smallmatrix}\alpha&\beta\\\gamma&\delta\end{smallmatrix}\right)$, the lower left
entries give $\gamma(1+\gamma)=0$. If $\gamma=0$ then $M$ is a power of $U$ and the image of
$\tilde q_{\Z_\ell}\circ\Phi_F$ is abelian, which is false since $\Phi_F$ is surjective. So $\gamma=-1$, the upper
left entries give $\beta=(\alpha-1)^2$, and $M=U^tLU^{-t}$ with $t=1-\alpha$. As $U^t$ commutes with $U$, this
proves the existence of $t_F$; it is unique because a power of $U$ commuting with $L$ is trivial; additivity
follows. Modulo $4$ the element $q_{12}(f)$ lies in the commutator subgroup of $V_4$, which is trivial; so
$M\equiv L$ and $t\equiv0$. Modulo $3$: apply $q$ to $\Phi_F(s)=\beta s\beta^{-1}$. In $\PSL_2(\F_3)=A_4$ the
element $U^{-t}q(\beta)$ centralizes the involution $q(s)$, so it lies in the four-group, the image of the
quaternion group $Q_8\subset\SL_2(\F_3)$. In $\SL_2(\F_3)/Q_8=\Z/3$ the classes of $U$ and $L=S_0US_0^{-1}$ are
equal, say to $1$; so $q_{12}(\beta)$, with $[\beta]=(-\rho_2,-\rho_2)$, has class $-4\rho_2=-\rho_2$, and
$-t-\rho_2=0$.
\end{proof}

\begin{definition}\label{def:delta}
$H_0=\{F\in\GTz:\lambda_F=1,\ t_F=0\}$. For $F\in H_0$ we have $q\circ\Phi_F=q$ on $\Bh$ and
$q_{12}\circ\Phi_F=q_{12}$ on $\Fh$ (Lemma~\ref{lem:E}(a)). So $q_{12}(A_F)$ and $q_{12}(\beta_F)$ lie in the groups
of Lemma~\ref{lem:E}(b),(c). Define $a_3(F),a_2(F)\in\F_2$ by
\[
 q_{12}(A_F)\bmod3\in(-I)^{a_3(F)}\langle R_0\rangle,\qquad
 q_{12}(\beta_F)\bmod3\in S_0^{a_2(F)}\langle-I\rangle,
\]
and $\delta=a_2+a_3$.
\end{definition}

By Lemma~\ref{lem:frames} and the invariance of $q_{12}$, the maps $F\mapsto q_{12}(A_F)$ and
$F\mapsto q_{12}(\beta_F)$ are homomorphisms from $H_0$ to abelian groups. So $a_2,a_3,\delta$ are continuous
characters of $H_0$, a closed subgroup of index $1$ or $3$ in $\{\lambda=1\}$. By Lemma~\ref{lem:shear},
$\Lambda_0\subset H_0$.

\subsection*{The lattice of the group $(2,3,6)$}
Let $\omega=e^{2\pi i/3}$, $\zeta=e^{2\pi i/6}$, $\mathcal L=\Z[\omega]$, and let $\mathcal L\rtimes\mu_6$ be the
group of the maps $(t,\xi):w\mapsto\xi w+t$ of $\mathbf C$. The assignment $u\mapsto(0,\zeta)$, $v\mapsto(1,\zeta)$
defines a surjection $\pi_6:B\to\mathcal L\rtimes\mu_6$ whose kernel is the second derived group $B''$; it maps
$B'$ onto $\mathcal L$. One computes $\pi_6(s)=(\zeta,-1)$, $\pi_6(x)=(0,\omega)$, $\pi_6(y)=(1+\zeta,\omega)$; so
$\pi_6(F)\subset(1-\omega)\mathcal L\rtimes\langle\omega\rangle$ and $\pi_6(F')\subset3\mathcal L$.

\begin{lemma}\label{lem:unit}
Let $F\in H_0$. Then $\Phi_F$ acts on $\Bh'^{\,\ab}=\widehat{\mathcal L}$ as multiplication by a unit $\nu_F$ of
$\Zh[\omega]$ with
\[
 \nu_F\equiv1\bmod3,\qquad \nu_F\equiv(-1)^{a_2(F)}\bmod4 .
\]
\end{lemma}

\begin{proof}
$\Phi_F$ preserves the closure of $B''$ and fixes $u$; the induced automorphism of
$\widehat{\mathcal L}\rtimes\mu_6$ fixes $(0,\zeta)$, so it is $(t,\xi)\mapsto(\nu t,\xi)$ for a unit $\nu$.
With $\pi_6(f)=(\phi,1)$, $\phi\in3\widehat{\mathcal L}$, the image of $\Phi_F(v)=f^{-1}vf$ is
$(1+(\zeta-1)\phi,\zeta)$; so $\nu=1+\omega\phi\equiv1\bmod3$. Let $\pi_6(\beta_F)=(b,\omega^j)$. Conjugating
$(\zeta,-1)$ by it gives $(2b+\omega^j\zeta,-1)$, which must be $(\nu\zeta,-1)$: $2b=(\nu-\omega^j)\zeta$. The
map $q\bmod3:\Bh\to\PSL_2(\F_3)=A_4$ kills $B''$, so it factors through a surjection
$\kappa:\widehat{\mathcal L}\rtimes\mu_6\to A_4$; it maps $\widehat{\mathcal L}$ onto the normal four-group of $A_4$, so it kills $2\widehat{\mathcal L}$, and
$\kappa(0,-1)=q(u^{3})=1$. As $2\widehat{\mathcal L}\rtimes\{\pm1\}$ has index $12$, it is the kernel of
$\kappa$. By definition of
$a_2$, $q(\beta_F)=q(s)^{a_2}$ modulo $3$, that is, $\kappa(b,\omega^j)=\kappa(\zeta,-1)^{a_2}=\kappa(a_2\zeta,1)$.
Hence $j=0$ and $b\equiv a_2\zeta\bmod2$, and $\nu=1+2b/\zeta\equiv1+2a_2\bmod4$.
\end{proof}

\section{Normalized elements and mapping class groups of genus one}\label{sec:genusone}

Let $f\in\Lambda$. For a surface $\Sigma$ with marked points and two pants decompositions $P,Q$, choose a path of
A-moves and S-moves from $P$ to $Q$, the $i$-th replacing a curve with twist $b_i$ by one with twist $c_i$, and let
\[
 W_f(P,Q)=\prod_i f(b_i^{\epsilon_i},c_i^{\epsilon_i})^{-1}\in\Gh(\Sigma)_{\rm pure},
\]
the product taken in the order of the path, with $\epsilon_i=1$ for an A-move and $2$ for an S-move. By
\cite[Theorem 2 and Lemma 6]{HLS} this is independent of the path: the cycles that generate the relations give
the unit element, by (I), (II), (III) and (R). Their automorphism is $F_P(b)=WbW^{-1}$ for a twist $b$ along a
curve of $Q$, with $W=W_f(P,Q)$ \cite[formula (3)]{HLS}. Clearly $W_f(P,Q)W_f(Q,Q')=W_f(P,Q')$ and
$W_f(gP,gQ)=gW_f(P,Q)g^{-1}$ for every mapping class $g$, pure or not.

\begin{lemma}\label{lem:mixed}
Let $\Gamma$ be the mapping class group of $\Sigma$, with all permutations of the marked points allowed.
\begin{itemize}
\item[(a)] $\alpha_P(g)=W_f(P,gP)\,g$ defines a homomorphism $\Gamma\to\widehat\Gamma$, which extends to an
automorphism $\alpha_P$ of $\widehat\Gamma$. It induces the identity on permutations, restricts to $F_P$ on the
pure subgroup, and $\alpha_Q=\Ad(W_f(P,Q)^{-1})\alpha_P$.
\item[(b)] Let $p$ be the map forgetting a marked point, defined on its stabilizer, with stable target. Then
$p\circ\alpha_P=\alpha_{pP}\circ p$, where $pP$ is the image of $P$.
\item[(c)] In genus one with $n\ge1$ marked points, let $\per:\widehat\Gamma\to\SL_2(\Zh)$ be the action on the
homology of the closed torus. There is $\bar\eta\in\GL_2(\Zh)$ with
$\per(\alpha_P(g))=\bar\eta\,\per(g)\,\bar\eta^{-1}$ for all $g$. Moreover $\alpha_P$, extended by \textup{(b)} to
the group with one more marked point, maps each loop around a puncture, in the fundamental group $\widehat\Pi$ of
the punctured torus, to a $\widehat\Pi$-conjugate of itself.
\end{itemize}
\end{lemma}

\begin{proof}
(a) The two properties of $W_f$ give $\alpha_P(g)\alpha_P(g')=\alpha_P(gg')$. A homomorphism to a profinite
group extends to the completion. On a twist $T$ along a curve of $Q$, $W_f(P,TP)=W_f(P,Q)\,TW_f(Q,P)T^{-1}$, so
$\alpha_P(T)=W_f(P,Q)TW_f(P,Q)^{-1}=F_P(T)$. The extension is an automorphism because it is one on the open pure
subgroup \cite[Theorem 4]{HLS} and the identity on the quotient. (b) The image of a pants decomposition, after
deleting inessential curves and identifying parallel ones, is a pants decomposition. A move not involving the
forgotten point keeps its type and its weight. An A-move in a four-holed sphere one of whose holes is the
forgotten point has both curves mapped to boundary curves of one pair of pants; their twists commute, the weight
dies because $f$ is a commutator word, and the two decompositions have the same image. An S-move is supported in a
one-holed torus, which survives since the target is stable. So $p(W_f(P,Q))=W_f(pP,pQ)$. (c) Add a marked
point $*$. The kernel of forgetting $*$ is $\Pi=\pi_1(\Sigma,*)$, a free group, and it has its full profinite
topology in the larger group, by (F). By (b), $\alpha_P$ lifts to an automorphism of the larger completed group
that preserves $\widehat\Pi$; let $\eta$ be its restriction. Then $\alpha_P(g)$ acts on $\widehat\Pi$ by
$\eta\circ g\circ\eta^{-1}$, up to inner automorphisms. The loop $C_i$ around the $i$-th puncture is, as a mapping
class, the twist along the curve surrounding that puncture and $*$; by (a), $\eta(C_i)=WC_iW^{-1}$ with $W$ pure.
Discrete pure classes map $C_i$ to $\Pi$-conjugates of itself; the set of $\widehat\Pi$-conjugates is closed, so
the same holds for $W$. Hence $\eta$ preserves the kernel of $\widehat\Pi\to H_1(T,\Zh)$, and $\bar\eta$ is the
induced map.
\end{proof}

\begin{lemma}\label{lem:lowgenus}
Let $F=(1,f)$ be normalized.
\textup{(a)} On $\Gh_{1,1}=\Bh_3/\overline{\langle(uv)^6\rangle}$, with $P$ the curve of $u$, $\alpha_P=\Phi_F$.
\textup{(b)} On $\Gh_{1,2}=\Bh_4/\overline{\langle\Omega\rangle}$, with $a,b,c$ the twists along a chain and
$P=\{a,c\}$, $\alpha_P=\Ad(g)\circ A_F$ with $g=f((ab)^3,c^2)$.
\end{lemma}

\begin{proof}
(a) The curve of $v$ is obtained by one S-move, of weight $f(u^2,v^2)^{-1}$. (b) Put $\omega_3=(ab)^3$ and
$e_1=\omega_3^2$, the twist along the boundary of a neighbourhood of the first two curves. The path
$\{a,c\}\to\{a,e_1\}\to\{b,e_1\}$ consists of an A-move and an S-move, so $\alpha_P$ fixes $a,c$ and conjugates $b$
by $W=f(c,e_1)^{-1}f(a^2,b^2)^{-1}$. In $B_4$ one has $(\omega_3c)^2=\Omega a^2=(c\omega_3)^2$: indeed
$\Omega=\omega_3\cdot cba^2bc$ and $\omega_3=ba^2ba^2$. So $u\mapsto c$, $v^2\mapsto\omega_3$ defines a homomorphism
$J\to B_4$, and the normalized relation (IV) gives $f(c,e_1)=f(c^2,\omega_3)$. On the other hand
$A_F(c)=g'^{-1}cg'$ with $g'=f(ba^2b,c^2)$; since $ba^2b=\omega_3a^{-2}$ and $a$ commutes with $\omega_3$ and $c$,
$g'=f(\omega_3,c^2)=g$, and $g$ commutes with $a$. Thus $g=f(c,e_1)^{-1}$, and $\Ad(g)A_F$ fixes $a$ and $c$ and
conjugates $b$ by $gf(a^2,b^2)^{-1}=W$.
\end{proof}

\section{The hexagonal self-similarity}\label{sec:H}

Let $q:\Gh_{1,2}\to\Gh_{1,1}$ forget the second point; on generators $(a,b,c)\mapsto(u,v,u)$. Its kernel is
$\widehat K$, where $K=\pi_1(E\setminus O)$ is free on $\xi_1=ca^{-1}$ and $\xi_2=abac^{-1}b^{-1}a^{-1}$, with its
full topology, by (F); on $K^{\ab}=\Z^2$ the elements $a,c$ act by $U$ and $b$ by $L$. Let
$G_3=q^{-1}(\langle uv\rangle)$; the image of $uv$ in $\Gh_{1,1}$ has order six.

\begin{proposition}\label{prop:H}
Let $F=(1,f)$ be normalized, and let $\tilde A=A_F(a^2,b^2)\in\Gh_{1,2}$ be the frame of Lemma~\ref{lem:frames}
evaluated at $a^2,b^2$. Then $\Psi=\Ad(\tilde A^{-1})\circ A_F$ preserves $G_3$, and there are an isomorphism
$\varphi:\Bh\to G_3$ with $\varphi(\overline{B'})=\widehat K$, induced by an isomorphism of discrete groups, and an
element $h\in G_3$ such that
\[
 \Psi\circ\varphi=\Ad(h)\circ\varphi\circ\Phi_F .
\]
\end{proposition}

\begin{proof}
\emph{Step 1: the points of order three.} Let $T=\mathbf R^2/\Z^2$ with the nine points of $T[3]$ marked,
$\Gamma=\Gamma_{1,[9]}$, $P_1=\Gamma_{1,9}$, and $H=T[3]\subset\Gamma$ acting by translations. Fix a pants
decomposition and let $\alpha=\alpha_P$ on $\widehat\Gamma$ (Lemma~\ref{lem:mixed}). By Lemma~\ref{lem:BZ}(a) with
$p=3$ there are $e\in\widehat{P_1}$ and $h\mapsto h'$ with $\alpha(h)=eh'e^{-1}$, $h'\in P_1h$. By
Lemma~\ref{lem:mixed}(c), $\per(h')$ is conjugate to $\per(h)=1$. The group $\{h'\}$ acts conformally on a complex
torus $T'$ with nine marked points \cite{K}, trivially on homology, hence by translations, and simply transitively
on the marked points; so, moving the origin, the point labelled $h$ is $\psi(h)$ for an isomorphism
$\psi:H\to T'[3]$, and $h'$ is the translation by $\psi(h)$, transported by the marking $m:T\to T'$. Choose an
oriented linear isomorphism $\ell:T\to T'$; then $\ell^{-1}\psi$ is given by a matrix $D\in\GL_2(\F_3)$. Let
$A\in\GL_2(\Z)$ lift $D$, with $\det A=\varepsilon=\pm1$. Then $g=m^{-1}\ell A$ fixes every marked point, has
orientation $\varepsilon$, and $h'=ghg^{-1}$.

We claim $\varepsilon=1$. Let $\widehat\Pi$ be the completed fundamental group of $T\setminus T[3]$ and
$\eta$ as in Lemma~\ref{lem:mixed}(c). The outer automorphism $\theta$ of $\widehat\Pi$ induced by $\eta^{-1}eg$
commutes with the outer action of each $h\in H$. The quotient of $T\setminus T[3]$ by $H$ is a once-punctured
torus, with fundamental group free on two generators $A_1,B_1$ and $\Pi$ the kernel of its abelianization modulo
$3$. Conjugation embeds its completion in $\Aut(\widehat\Pi)$ as the preimage of the outer action of $H$, so
$\theta$ extends to an automorphism $\Theta$ of this free profinite group of rank two, inducing the identity on
$H$. A loop around the puncture lies in $\Pi$ and is one of the nine puncture loops of $T\setminus T[3]$; by
Lemma~\ref{lem:mixed}(c) and the orientation of $g$, $\Theta$ maps it to a conjugate of its $\varepsilon$-th
power. In the quotient of class two this loop is $[A_1,B_1]$, which $\Theta$ raises to the power $\det\Theta^{\ab}$.
So $\varepsilon=\det\Theta^{\ab}\equiv1\bmod3$, and $\varepsilon=1$.

Hence $c_1=(eg)^{-1}$ is pure and $\alpha'=\Ad(c_1)\circ\alpha$ fixes $H$ pointwise. The normalizer of $H$ in
$\Gamma$ is $H\rtimes\SL_2(\Z)$, the affine maps; its subgroup $\mathcal N_0=\SL_2(\Z)$ fixing the origin is mapped
isomorphically to $\Gamma_{1,1}$ by forgetting the eight other points, so its closure $\overline{\mathcal N_0}$ is
$\Gh_{1,1}$. By Lemma~\ref{lem:BZ}(b), $\alpha'$ preserves $N_{\widehat\Gamma}(H)=H\overline{\mathcal N_0}$, and it
preserves the stabilizer of the origin; so it preserves $\overline{\mathcal N_0}$. By Lemmas
\ref{lem:mixed}(b) and \ref{lem:lowgenus}(a),
\[
 \alpha'|_{\overline{\mathcal N_0}}=\Ad(c_2)\circ\Phi_F\quad\text{on }\Gh_{1,1},\qquad c_2\in\Gh_{1,1}.
\]
In particular $\alpha'$ fixes the central involution $\iota=-1$ of $\overline{\mathcal N_0}$.

\emph{Step 2: the quotient by $\iota$.} Let $C=C_\Gamma(\iota)$. The quotient $S=T/\iota$ is a sphere with four
branch points $b_0,\dots,b_3$, the images of $T[2]$, with $b_0$ the image of the origin, and four more marked
points $t_1,\dots,t_4$, the images of the pairs $\pm t$ of points of order three. By Birman and Hilden \cite{BH},
$C/\langle\iota\rangle$ is the group $\mathcal D$ of mapping classes of $S$ that fix $b_0$ and preserve
$\{b_1,b_2,b_3\}$ and $\{t_1,\dots,t_4\}$. We claim
\[
 C_{\widehat\Gamma}(\iota)=\overline C\cong\widehat C,\qquad \widehat C/\langle\iota\rangle=\widehat{\mathcal D}.
\]
The first equality is Lemma~\ref{lem:BZ}(b). For the second, mark also the three points of
$T[2]\setminus0$: let $\Gamma^+=\Gamma_{1,[9]+[3]}$ and $C^+=C_{\Gamma^+}(\iota)$. Forgetting the three points
maps $C^+$ isomorphically to $C$, since a homeomorphism commuting with $\iota$ preserves its fixed points
\cite{BH}. Now
$T$ minus the twelve points is an unramified double cover of $S$ minus eight points; let $\Pi^+\subset\Phi$ be
their fundamental groups. The completion of $\Gamma_{0,[8]}$ acts faithfully on $\widehat\Phi$ by outer
automorphisms: for the pure group this is Asada's theorem in genus zero \cite{As}, and an element with a nontrivial
permutation moves the conjugacy classes of the puncture loops. Hence so does $\widehat{\mathcal D}$. If $z\in\widehat{C^+}$ acts trivially on
$\widehat{\Pi^+}$, its image in $\widehat{\mathcal D}$ is represented by an automorphism of $\widehat\Phi$ that is
the identity on the open subgroup $\widehat{\Pi^+}$; such an automorphism is the identity, because the centralizer
of an open subgroup of a nonabelian free profinite group is trivial. So $z\in\langle\iota\rangle$, and $z=1$
because $\iota$ permutes the punctures. Thus $\widehat{C^+}$ embeds in $\widehat{\Gamma^+}$. Finally the kernel
$\widehat{K_3}$ of $\widehat{\Gamma^+}\to\widehat\Gamma$ is the completion of the fundamental group $K_3$ of the
space of three unordered points of $T\setminus T[3]$, by (F) applied to the three successive point-pushing
kernels for ordered points and passage to a subgroup of finite index; the discrete centralizer of $\iota$ in $K_3$ is
trivial, so by Lemma~\ref{lem:BZ}(b), applied to $K_3\rtimes\langle\iota\rangle$, its centralizer in
$\widehat{K_3}$ is trivial. Hence $\widehat{C^+}\to\overline C$ is injective, which proves the claim; the last
equality is right exactness of completion.

So $\alpha'$ induces an automorphism $\delta_8$ of $\widehat{\mathcal D}$. The pure subgroup $\Gamma_{0,8}$ of
$\mathcal D$ is the image of the classes in $C$ that fix each pair $\pm t$ and act trivially on $T[2]$. By
Lemma~\ref{lem:mixed}(a),(c), $\alpha'$ preserves the closure of this group, and $\delta_8$ acts on
$\mathcal D/\Gamma_{0,8}=\mathfrak S_{\{b_1,b_2,b_3\}}\times\mathfrak S_{\{t_j\}}$ by conjugation by a permutation
$\sigma_2$ of $b_1,b_2,b_3$.

\emph{Step 3: the parameter.} By Hoshi, Minamide and Mochizuki
\cite[Corollary C; Corollary 2.8 in the body]{HMM}, for the curve of type $(g,r)=(0,3)$ and $n\ge2$ the outer
automorphism group of the profinite group $\Pi_n=\Gh_{0,n+3}$ is $\GT\times\mathfrak S_{n+3}$. With $n=5$:
$\Out(\Gh_{0,8})=\GT\times\mathfrak S_8$, the first factor acting by $A_{F'}$ and the second by permutations of the
marked points. That this decomposition is compatible with forgetting points is not part of their statement; it
follows from (B3). Let $(F',\sigma)$ be the class of
$\delta_8|\Gh_{0,8}$. Conjugation by an element of $\widehat{\mathcal D}$ with permutation $\pi$ has class $(1,\pi)$;
comparing with its image under $\delta_8$ gives $\sigma\pi\sigma^{-1}=\sigma_2\pi\sigma_2^{-1}$ for all $\pi$ in
$\mathfrak S_3\times\mathfrak S_4$, whose centralizer in $\mathfrak S_8$ is trivial. So $\sigma=\sigma_2$. To find
$F'$, fix $j$ and forget the points $t_i$, $i\ne j$. On the torus this is forgetting all marked points except the
pair $\pm t_j$, which lands in $\Gamma_{1,[2]}=\Gamma_{1,2}\times\langle\iota\rangle$; and
$\Gamma_{1,[2]}/\langle\iota\rangle=\Gamma_{1,2}=B_4/\langle\Omega\rangle$ is the group of classes of $S$ fixing
$t_j$ and preserving the set of branch points, with pure subgroup $\Gamma_{0,5}$ \cite{BH}. By Lemmas
\ref{lem:mixed}(b) and \ref{lem:lowgenus}(b), $\alpha'$ induces on $\Gh_{1,2}$ the automorphism $A_F$ up to an
inner automorphism; so the class of the induced automorphism of $\Gh_{0,5}$ is $(F,\pi)$ for some $\pi$. It is also
$(F',\sigma_2)$. So $F'=F$.

\emph{Step 4: forgetting the branch points and taking the double cover.} Let $\mathcal D_5$ be the group of
classes of $S$ fixing $b_0$ and preserving $\{t_1,\dots,t_4\}$; it is $B_4/\langle\Omega\rangle$, the braids of
the $t_j$ in the disc $S\setminus b_0$, and its pure subgroup is $\Gamma_{0,5}$. By Step 3, $\delta_8$ preserves the
kernel of $\Gh_{0,8}\to\Gh_{0,5}$ and induces there the class $(F,1)$. An element of $\widehat{\mathcal D_5}$
that centralizes $\Gh_{0,5}$ is trivial: its permutation is trivial by \cite{HMM}, and $\Gh_{0,5}$ has trivial
centre. It follows that $\delta_8$ maps the kernel of $\widehat{\mathcal D}\to\widehat{\mathcal D_5}$ into itself
(the image of an element $k$ of the kernel commutes with the image of $\Gh_{0,8}$), that it induces an automorphism
$\delta_5$ of $\widehat{\mathcal D_5}$, and that $\delta_5=\Ad(c_3)\circ A_F$: two automorphisms that agree on
$\Gh_{0,5}$ agree. Let $E'\to S$ be the double cover branched at $t_1,\dots,t_4$, and $O',P'$ the points above
$b_0$. By \cite{BH}, lifting classes so as to fix $O'$ and $P'$ identifies $\mathcal D_5$ with
$\Gamma_{1,2}$, the generators of $B_4$ becoming the twists $a,b,c$ along a chain.

\emph{Step 5: the composite map.} The composite
$\mathcal N_0\subset C\to\mathcal D\to\mathcal D_5=\Gamma_{1,2}$ kills $\iota$ and defines $\varphi:B\to\Gamma_{1,2}$,
and by construction $\delta_5\circ\varphi=\varphi\circ(\alpha'|_{\overline{\mathcal N_0}})$ on $\Bh$. It is the map
on orbifold fundamental groups induced by the map of moduli $(E,O)\mapsto(E';O',P')$. For
$E:y^2=x^3+a_4x+a_6$, the points $t_j$ are the roots of the $3$-division polynomial
$3x^4+6a_4x^2+12a_6x-a_4^2$, the curve $E'$ is $t^2=x^4+2a_4x^2+4a_6x-a_4^2/3$, and with
$X=(t+x^2+a_4)/2$, $Z=X-a_4/3$, $Y=Xx+a_6/2$ one finds
\[
 Y^2=Z^3+\tfrac1{108}(4a_4^3+27a_6^2),\qquad O'=\infty,\quad P'=(-a_4/3,\,-a_6/2).
\]
So $E'$ is always the curve $E_0:Y^2=Z^3+1$ up to isomorphism, and $(E,O)\mapsto P'$ identifies the moduli of
elliptic curves up to $\pm1$ with the quotient of $E_0\setminus O'$ by $\mu_6=\Aut(E_0)$. This quotient is the fibre
of the moduli of two-pointed curves over the point $E_0$. Therefore $\varphi$ is injective, its image is the
preimage under $q$ of a cyclic subgroup of order six of $\Gamma_{1,1}$, which after conjugation is
$\langle uv\rangle$, and $\varphi(B')=K$. Passing to completions, $\varphi:\Bh\to G_3$ is an isomorphism.

\emph{Step 6: the conjugator.} By Steps 1, 4 and 5, $A_F\circ\varphi=\Ad(h_0)\circ\varphi\circ\Phi_F$ for some
$h_0\in\Gh_{1,2}$. In $\Bh_3$ we have $\Phi_F(uv)=\tilde A_3(uv)\tilde A_3^{-1}$ with $\tilde A_3=A_F(u^2,v^2)$,
because both sides agree modulo the centre and in $\Bh_3^{\ab}$. The restriction of $A_F$ to
$\overline{\langle a,b\rangle}=\Bh_3$ is $\Phi_F$, so $\Psi$ fixes $ab$. Also $q\circ A_F=\Phi_F\circ q$, since
$q(f(ba^2b,c^2))=f((uv)^3u^{-2},u^2)=1$. So $\Psi$ preserves $\widehat K$ and $G_3$, and
$\Psi\varphi=\Ad(h)\varphi\Phi_F$ with $h=\tilde A^{-1}h_0$. Apply $q$: since $\Phi_F$ fixes $\tau$, the element
$q(h)$ centralizes $q(\varphi(\Bh))=\langle uv\rangle$. Modulo the central involution
$(uv)^3$, the group $\Gh_{1,1}$ is $\Bh=\widehat{C_2}\amalg\widehat{C_3}$ and $uv$ generates the factor of order
three, whose centralizer is that factor \cite{HR}; the preimage of this factor in $\Gh_{1,1}$ is
$\langle uv\rangle$. So $q(h)\in\langle uv\rangle$ and $h\in G_3$.
\end{proof}

\section{Proof of Theorem~\ref{thm:delta}}\label{sec:proof}

Let $F$ be normalized; then $F\in H_0$. We compute the action of the two sides of Proposition~\ref{prop:H} on
abelianizations modulo $12$.

\emph{The right side.} $\varphi$ induces an isomorphism $C:\widehat{\mathcal L}=\Bh'^{\,\ab}\to\widehat K^{\ab}$, and
$C\zeta C^{-1}$ is the action of $\varphi(\tau)$, which is $R_+^{\pm1}$. By Lemma~\ref{lem:unit},
$\varphi\Phi_F\varphi^{-1}$ acts on $\widehat K^{\ab}/12$ by $C\nu_FC^{-1}$, which is $(-1)^{a_2(F)}I$ modulo $4$ and
$I$ modulo $3$. The element $h\in G_3$ acts by $R_+^k$ for some $k\in\Z/6$.

\emph{The left side.} For every $g\in\Gh_{1,2}$ the action of $A_F(g)$ on $\widehat K^{\ab}/12$ is
$\tilde q(\Phi_F(q(g)))=\tilde q(q(g))$, because $t_F=0$. So the action $N$ of $A_F$ on $\widehat K^{\ab}/12$
commutes with $\SL_2(\Z/12)$ and is a scalar. Since $A_F(ca^{-1})=g^{-1}(ca^{-1})g$ with
$g=f(ba^2b,c^2)\in\widehat K$, $N$ fixes the class of $\xi_1$, and $N=I$. The element $\tilde A$ acts by
$A_F(U^2,L^2)=q_{12}(A_F)$. So $\Psi$ acts by $q_{12}(A_F)^{-1}$, which is $I$ modulo $4$ and lies in
$(-I)^{a_3(F)}\langle R_0\rangle$ modulo $3$ (Lemma~\ref{lem:E}(b)).

\emph{Comparison.} Modulo $4$: $I=R_+^k\cdot(-1)^{a_2}$. As $R_+$ has order six modulo $4$ and $R_+^3=-I$, this
gives $k\equiv3a_2\bmod6$. Modulo $3$: $q_{12}(A_F)^{-1}=R_+^k=(-1)^kR_0^k$. On the line fixed by $R_0$ the left
side acts by $(-1)^{a_3}$ and the right side by $(-1)^k=(-1)^{a_2}$. So $a_2(F)=a_3(F)$. \qed

\section{What remains open}

\begin{enumerate}
\item Whether $(\mathrm{III}')$ holds on $\GT$. With Theorem~\ref{thm:IV} this is the question whether
$\mathrm{I}\Gamma=\GT$. The method of Section~\ref{sec:IV} applies to any locus of curves with automorphisms seen
through two forgetting maps; for $(\mathrm{III}')$ the natural loci are those of order two and five in the moduli
of five points, and we have not found the two maps.
\item Whether the action of $\Lambda_0$ and the Galois action fit into one action of $\GT$ on the tower with
marked points \cite[\S8]{LNS}. The obstruction is that the decomposition $F=F_*\sigma$ is not unique and that the
Galois action is not given by the formulas of \cite{HLS} when $\rho_2\ne0$.
\item The statements quoted in Section 1 about the character on all of $\GTz$: its extension, its vanishing on
$\GQ$ and its nonvanishing. They need their own check.
\end{enumerate}

\section*{Provenance and checks}

The mathematics is the work of GPT-6 Astra (ChatGPT Pro mode), in research runs of November 2026 on notes of
Grothendieck's \emph{La longue marche \`a travers la th\'eorie de Galois}. The runs contain two proofs of
Theorem~\ref{thm:delta}; this note follows the one through $\Lambda$.

Claude Fable 5.2 at effort max checked every step of the runs' proofs of Theorems \ref{thm:IV} and \ref{thm:Lambda}
(18 November), re-proved the computation of Section~\ref{sec:proof} independently (21 November), and checked the
chain behind Proposition~\ref{prop:H}, re-deriving Lemma~\ref{lem:lowgenus}(b) (23 November). Claude Opus 5.5 at
effort max checked the parity computation, exhaustively over $\Z/12$ with its own program, and recomputed the
identity of Step 5 (22 November). As second reader of this text (23 November) it followed every step of Sections
4 and 8 and of the proofs rewritten by the editor, re-derived the matrix and braid identities, opened
\cite{HMM} at the source, and found no mathematical error; its corrections of citation and precision are
incorporated. No fault was found in the statements. Before deposit Claude Opus 5.5 at effort max (23 November)
read the text, recomputed the abelianizations and the character argument of Propositions~\ref{prop:IVc} and
\ref{prop:c}, and checked the Galois decomposition in the proof of Theorem~\ref{thm:Lambda}.

This text was written by Claude (Claude Fable 5.2 at effort max) as editor. It differs from the runs' texts in
these places, which the second reader checked: the finite-subgroup steps cite \cite{BZ} (Lemma~\ref{lem:BZ}) where one
of the runs proves them by Smith theory; the shear $t_F$ is defined by Lemma~\ref{lem:shear}; Lemma~\ref{lem:unit}
and the first half of the proof of Proposition~\ref{prop:c} are new derivations; in Step 2 the twelve points are
used only for the topology of the centralizer.

Read at the source for this text: \cite{NS} (\S\S1, 5, 6, author's version), \cite{HLS} (pp.~3--5, 18--19),
\cite{S} (pp.~659, 670), \cite{LNS} (\S\S1, 2, 8), \cite{BZ} (statements and the proof of Theorem A). \cite{HMM} was
opened by the second reader (Corollary C and Corollary 2.8). Of \cite{BBR}, \cite{DG} and \cite{Co} only the
abstracts were read; \cite{BBR} also in its first section. The bibliographic data of \cite{As} and \cite{FR} were
confirmed; \cite{A} was not opened, and what is taken from it is proved in (F). The other references are cited
from memory or with the locators of the runs. No human has checked the proofs or this
text.

The records of the checks are at the permanent page of this work, WB-0007:

{\raggedright\smallskip\noindent\url{\recordurl}.\par}

\smallskip
This text is dedicated to the public domain (CC0 1.0).

\begin{thebibliography}{99}\small
\bibitem{A} M. P. Anderson, Exactness properties of profinite completion functors, \emph{Topology} 13 (1974),
229--239.
\bibitem{As} M. Asada, The faithfulness of the monodromy representations associated with certain families of
algebraic curves, \emph{J. Pure Appl. Algebra} 159 (2001), no.~2, 123--147.
\bibitem{BBR} L. Basualdo Bonatto and M. Robertson, Galois actions on surfaces and a higher genus
Grothendieck--Teichm\"uller group, arXiv:2606.01466v2 (2026).
\bibitem{BH} J. S. Birman and H. M. Hilden, On isotopies of homeomorphisms of Riemann surfaces, \emph{Ann. of Math.}
97 (1973), 424--439.
\bibitem{BZ} M. Boggi and P. Zalesskii, Finite subgroups of the profinite completion of good groups, \emph{Bull.
London Math. Soc.} 57 (2025), 236--255; arXiv:2406.08639.
\bibitem{Co} B. Collas, Action of a Grothendieck--Teichm\"uller group on torsion elements of full Teichm\"uller
modular groups of genus zero, \emph{J. Th\'eor. Nombres Bordeaux} 24 (2012), 605--622.
\bibitem{DG} V. A. Dolgushev and J. J. Guynee, GT-shadows for the gentle version of the Grothendieck--Teichm\"uller
group, arXiv:2401.06870 (2024).
\bibitem{D} V. G. Drinfeld, On quasitriangular quasi-Hopf algebras and on a group that is closely connected with
$\mathrm{Gal}(\overline{\mathbf Q}/\mathbf Q)$, \emph{Leningrad Math. J.} 2 (1991), 829--860.
\bibitem{FR} D. I. Fouxe-Rabinovitch, \"Uber die Automorphismengruppen der freien Produkte I, \emph{Mat. Sb.} 8(50)
(1940), 265--276.
\bibitem{HS} D. Harbater and L. Schneps, Fundamental groups of moduli and the Grothendieck--Teichm\"uller group,
\emph{Trans. Amer. Math. Soc.} 352 (2000), 3117--3148.
\bibitem{HLS} A. Hatcher, P. Lochak and L. Schneps, On the Teichm\"uller tower of mapping class groups, \emph{J.
reine angew. Math.} 521 (2000), 1--24.
\bibitem{HR} W. Herfort and L. Ribes, Torsion elements and centralizers in free products of profinite groups,
\emph{J. reine angew. Math.} 358 (1985), 155--161.
\bibitem{HMM} Y. Hoshi, A. Minamide and S. Mochizuki, Group-theoreticity of numerical invariants and distinguished
subgroups of configuration space groups, \emph{Kodai Math. J.} 45 (2022), 295--348. Corollary C of the
introduction is Corollary 2.8 of the body.
\bibitem{I} Y. Ihara, On the embedding of $\mathrm{Gal}(\overline{\mathbf Q}/\mathbf Q)$ into $\GT$, in \emph{The
Grothendieck theory of dessins d'enfants}, London Math. Soc. Lecture Note Ser. 200, Cambridge Univ. Press, 1994,
289--321.
\bibitem{K} S. P. Kerckhoff, The Nielsen realization problem, \emph{Ann. of Math.} 117 (1983), 235--265.
\bibitem{LNS} P. Lochak, H. Nakamura and L. Schneps, The Grothendieck--Teichm\"uller group $\GT$ acts on the genus
$g$ mapping class group with $0$ or $1$ marked point, arXiv:2602.12462 (2026).
\bibitem{LS} P. Lochak and L. Schneps, The Grothendieck--Teichm\"uller group and automorphisms of braid groups, in
\emph{The Grothendieck theory of dessins d'enfants}, London Math. Soc. Lecture Note Ser. 200, Cambridge Univ.
Press, 1994, 323--358.
\bibitem{NS} H. Nakamura and L. Schneps, On a subgroup of the Grothendieck--Teichm\"uller group acting on the
tower of profinite Teichm\"uller modular groups, \emph{Invent. Math.} 141 (2000), 503--560.
\bibitem{S} L. Schneps, Automorphisms of curves and their role in Grothendieck--Teichm\"uller theory, \emph{Math.
Nachr.} 279 (2006), 656--671.
\end{thebibliography}

\end{document}
